\vspace*{6cm}
\begin{center}
{\Large \gesperrt{PREFACE.}}

\vspace{1em}
\rule{2cm}{0.4pt}
\end{center}
\vspace{2em}

It has been granted to me, after manifold hindrances, to carry out after all the plan for a continuation of my Textbook of Algebra that I announced twelve years ago in the preface to the first edition of the second volume. Thanks to the obliging attitude of the publishing firm, this third volume of the Algebra appears at the same time as the second edition of the ``Elliptic Functions and Algebraic Numbers'', first printed in the year 1891.

It is chiefly concerned with the further development of the manifold applications of algebra, and in particular of the theory of quadratic fields, to the problems that have arisen from the elliptic functions. These problems furnish us with the first example, going beyond cyclotomy, of algebraic numbers whose laws are known to some extent. A more thorough treatment of the quadratic fields with negative discriminant serves as the foundation for this. Admittedly, here too not everything has been achieved that I had set myself as the ultimate goal. Thus the elaboration of the theory of relatively cyclic fields still had to be postponed --- only for the time being, I hope.

On the other hand, in accordance with a wish that has repeatedly been put to me, I have added an outline of the theory of algebraic functions on an arithmetical foundation, which
